% Figure: Rankine cycle on the T-s diagram, with the water/steam saturation
% dome drawn schematically. State 1 saturated liquid, 2 compressed liquid,
% 3 superheated vapour, 4 wet vapour after the turbine.
\begin{figure}[htbp]
\centering
\begin{tikzpicture}
\begin{axis}[
  width=11cm, height=8cm,
  xlabel={specific entropy $s$ (kJ/kg$\cdot$K)},
  ylabel={temperature $T$ ($^{\circ}$C)},
  xmin=0, xmax=9, ymin=0, ymax=620,
  axis lines=left,
  xtick=\empty, ytick={33,264,450},
  yticklabels={$33$,$264$,$450$},
  tick label style={font=\scriptsize},
  label style={font=\small},
  clip=false,
]
% saturation dome: critical point near (4.4, 374). Left branch (sat. liquid),
% right branch (sat. vapour).
\addplot[enggray, thick, smooth] coordinates {
  (0,0) (0.5,40) (1.3,120) (2.3,230) (3.3,320) (4.0,365) (4.4,374)
};
\addplot[enggray, thick, smooth] coordinates {
  (4.4,374) (5.0,360) (6.0,300) (7.0,200) (7.8,100) (8.4,33)
};
\node[enggray, font=\scriptsize, anchor=south] at (axis cs:4.4,378) {critical pt};
% Rankine cycle path:
% 1 saturated liquid at condenser pressure (T=33C, s low ~0.47)
% 2 compressed liquid after pump (almost same s, slightly higher T)
% boiler 2->3: heat to saturated liquid, through dome, superheat to 3
% 3 superheated vapour at 450C, s ~ 6.8
% turbine 3->4 isentropic down to condenser pressure, s constant ~6.8, T=33
% condenser 4->1 along T=33 isotherm back to sat. liquid
\addplot[engblue, very thick] coordinates {(0.47,33) (0.5,38)};         % pump 1->2
\addplot[engblue, very thick, smooth] coordinates {                    % boiler 2->3
  (0.5,38) (1.6,264) (3.2,264) (5.0,360) (6.6,450)
};
\addplot[engblue, very thick] coordinates {(6.6,450) (6.6,33)};        % turbine 3->4 (isentropic)
\addplot[engblue, very thick] coordinates {(6.6,33) (0.47,33)};        % condenser 4->1
% shade the work area faintly
% state markers
\filldraw[engred] (axis cs:0.47,33) circle (1.6pt) node[anchor=north east, font=\scriptsize]{1};
\filldraw[engred] (axis cs:0.5,38)  circle (1.6pt) node[anchor=south east, font=\scriptsize]{2};
\filldraw[engred] (axis cs:6.6,450) circle (1.6pt) node[anchor=south west, font=\scriptsize]{3};
\filldraw[engred] (axis cs:6.6,33)  circle (1.6pt) node[anchor=north west, font=\scriptsize]{4};
\node[font=\scriptsize, engblue, anchor=south] at (axis cs:3.3,275) {boiler isobar};
\node[font=\scriptsize, engblue, anchor=south] at (axis cs:3.4,30) {condenser isobar};
\end{axis}
\end{tikzpicture}
\caption[Rankine cycle on the temperature-entropy diagram]{The ideal Rankine
cycle traced over the water saturation dome (gray). The pump raises the
saturated liquid from condenser to boiler pressure (1 to 2, a near-vertical
sliver). The boiler heats the water to saturation, evaporates it across the
dome, and superheats the vapour (2 to 3). The turbine expands the vapour
isentropically (3 to 4, vertical because entropy is constant). The condenser
returns the wet vapour to saturated liquid along the low isotherm (4 to 1).
Raising the boiler pressure or temperature lifts the average temperature of
heat addition; lowering the condenser pressure drops the temperature of heat
rejection. Both raise efficiency.}
\label{fig:vol04ch03:rankine-ts}
\end{figure}
